\chapter{Probabilistic Models and Uncertainty}\label{ch:ml:probabilistic-models-and-uncertainty}

Two forecasts both say a stock will rise five basis points tomorrow. One comes with a standard deviation of 80 basis
points, the other with 250: the same expected return on a quiet day and on a violent one. A book that sizes the two
trades alike holds three times the risk per unit of expected return in the second, and a risk report that shows only
the point forecast cannot tell them apart. Every forecast of a return is a distribution. This chapter trains models
that say so, scores them with losses that reward honest distributions, builds intervals whose coverage survives a
change of regime, and measures what the variance is worth when positions are sized.

\section{Predictive distributions and proper scoring rules}

\begin{definition}[Predictive distribution, proper scoring rule]\label{def:ml:probabilistic-models-and-uncertainty:scoring}
A \emph{predictive distribution}\index{predictive distribution} is a forecast of the whole conditional distribution
$F(\cdot\mid x)$ of the target. A scoring rule $S(F, y)$ gives a loss to a forecast $F$ when $y$ occurs; it is a
\emph{proper scoring rule}\index{proper scoring rule} if its expected value under the true distribution $G$ is minimised
by forecasting $F = G$, so that honesty is optimal (Gneiting and Raftery, 2007). The negative log-likelihood is proper;
so are the two rules below.
\end{definition}

\begin{definition}[Pinball loss, continuous ranked probability score]\label{def:ml:probabilistic-models-and-uncertainty:crps}
The \emph{pinball loss}\index{pinball loss} of a forecast $q$ of the $\alpha$-quantile is
$\ell_\alpha(y, q) = \max\{\alpha(y - q), (\alpha - 1)(y - q)\}$. The \emph{continuous ranked probability
score}\index{continuous ranked probability score} of a forecast distribution $F$ is
$\mathrm{CRPS}(F, y) = \int(F(z) - \mathbf 1\{z\ge y\})^2\,dz = 2\int_0^1\ell_\alpha(y, F^{-1}(\alpha))\,d\alpha$, in the
units of the target; for a Gaussian forecast it has a closed form.
\end{definition}

\begin{proposition}[The pinball loss elicits the quantile]\label{prop:ml:uncert:pinball}
If $Y$ has a continuous distribution function $G$, $\E\,\ell_\alpha(Y, q)$ is minimised at $q = G^{-1}(\alpha)$.
\end{proposition}

\begin{proof}
$\frac{d}{dq}\E\,\ell_\alpha(Y, q) = -\alpha\P(Y > q) + (1 - \alpha)\P(Y < q) = G(q) - \alpha$, which is zero at
$G^{-1}(\alpha)$, negative before and positive after.
\end{proof}

The chapter's data are \texttt{firm.mlsynth}'s daily series: twenty assets, ten years, GARCH volatility, Student-$t$
shocks with five degrees of freedom, a planted drift, and a regime from day 2\,100 in which every shock is 1.6 times
larger. The target is the next day's return; the features are the drift reading, the 5- and 20-day returns, the ratio of
short to long realised volatility and an EWMA volatility, all known at the close. Training days 60--1\,511, calibration
days 1\,512--1\,763, test days 1\,764--2\,518: the regime changes about half-way through the test.

\section{Quantile and distributional losses}

\begin{definition}[Quantile regression, mixture density network]\label{def:ml:probabilistic-models-and-uncertainty:quantile}
\emph{Quantile regression}\index{quantile regression} fits the conditional $\alpha$-quantile by minimising the average
\omterm{def:ml:probabilistic-models-and-uncertainty:crps}{pinball loss} $\mathcal L_n = \frac1n\sum_i\ell_\alpha(y_i, f_\theta(x_i))$ (Koenker and Bassett, 1978); with several
levels it describes the whole distribution. A \emph{mixture density network}\index{mixture density network} outputs the
weights, means and standard deviations of a mixture of $K$ Gaussians and is trained by their negative log-likelihood
(Bishop, 1994).
\end{definition}

Three distributional models are fitted: gradient-boosted trees for nineteen quantile levels (5\% to 95\%), a network with a
mean and a log-variance output trained by the Gaussian negative log-likelihood (five seeds, averaged as a \omterm{def:ml:neural-networks-for-noisy-tabular-data:ensemble}{deep ensemble},
chapter~7), and a two-component \omterm{def:ml:probabilistic-models-and-uncertainty:quantile}{mixture density network}. \Cref{tab:ml:uncert:scores} scores them with the CRPS on the
nineteen quantiles and the coverage of their 90\% intervals, before and after the regime change, beside the true
conditional distribution's Gaussian approximation.

\begin{table}[htbp]
\centering\small
\begin{tabular}{@{}lcccc@{}}
\toprule
 & \multicolumn{2}{c}{before day 2\,100} & \multicolumn{2}{c}{after day 2\,100}\\
\cmidrule(lr){2-3}\cmidrule(l){4-5}
model & CRPS (bp) & 90\% coverage & CRPS (bp) & 90\% coverage\\
\midrule
boosted quantiles & 78.9 & 89.4\% & 136.3 & 85.0\%\\
Gaussian network (one seed) & 79.0 & 90.6\% & 135.3 & 87.0\%\\
ensemble of five & 79.1 & 91.1\% & 135.4 & 87.6\%\\
\omterm{def:ml:probabilistic-models-and-uncertainty:quantile}{mixture density network} & 78.8 & 89.7\% & 136.1 & 86.1\%\\
\midrule
truth (Gaussian approximation) & 78.9 & 91.2\% & 135.5 & 91.3\%\\
\bottomrule
\end{tabular}
\caption{Distributional forecasts of the next day's return on the test days of twenty synthetic assets. CRPS computed on
nineteen quantiles (5\% to 95\%), in basis points of return. Data: \texttt{ml\_uncert.scores}.}\label{tab:ml:uncert:scores}
\end{table}

Before the break the four models are as good as the truth's approximation, within a quarter of a basis point of
CRPS: the volatility they must forecast is persistent and the EWMA feature carries most of it. After the break the CRPS
of all of them nearly doubles, as the truth's does, but their intervals cover 85--88\% of the days instead of 90\%: the
features adapt with a lag and the models were fitted on a calmer world. The mean forecasts are another matter; a
signal-to-noise ratio of a few hundredths hardly moves any of these scores, which are dominated by the volatility.

\section{Two kinds of uncertainty}

\begin{definition}[Aleatoric and epistemic uncertainty]\label{def:ml:probabilistic-models-and-uncertainty:aleatoric}
\emph{Aleatoric uncertainty}\index{aleatoric uncertainty} is the randomness of the target given the inputs, which more
data cannot remove; \emph{epistemic uncertainty}\index{epistemic uncertainty} is the model's uncertainty about the
function itself, which more data reduce. In a \omterm{def:ml:neural-networks-for-noisy-tabular-data:ensemble}{deep ensemble} the average of the members' predicted variances estimates
the first and the variance of their means the second; the predictive variance is their sum.
\end{definition}

On the chapter's test days the epistemic part is 0.09\% of the predictive variance. Daily returns are almost all
aleatoric: nothing a model learns will narrow them much. The epistemic part matters where the model is extrapolating
(inputs outside the training range, a new instrument) and in the few-label settings of chapter~10; it is the part an
ensemble, not a single network, can report (Kendall and Gal, 2017).

\section{Conformal prediction}

\begin{definition}[Conformal prediction]\label{def:ml:probabilistic-models-and-uncertainty:conformal}
\emph{Conformal prediction}\index{conformal prediction} turns any forecast into prediction sets with a coverage
guarantee: a score $s_i = s(x_i, y_i)$ (for example $|y_i - \hat y_i|$, or that divided by a predicted standard
deviation) is computed on $n$ calibration examples, and the set for a new $x$ is $\{y: s(x, y)\le\hat q\}$ with $\hat q$ the
$\lceil(n+1)(1 - \alpha)\rceil$-th smallest calibration score (split conformal; Vovk, Gammerman and Shafer, 2005). The
adaptive version updates the level online, $\alpha_{t+1} = \alpha_t + \gamma(\alpha - \mathrm{err}_t)$, to keep the
long-run error rate at $\alpha$ when the data drift (Gibbs and Candès, 2021).
\end{definition}

\begin{theorem}[Split conformal coverage]\label{thm:ml:uncert:conformal}
If the calibration scores and the new score are exchangeable, $\P(s_{n+1}\le\hat q)\ge1 - \alpha$.
\end{theorem}

\begin{proof}
By exchangeability the rank of $s_{n+1}$ among the $n + 1$ scores is uniform on $\{1,\dots,n+1\}$ (ties broken at random).
$s_{n+1}\le\hat q$ whenever its rank is at most $k = \lceil(n+1)(1-\alpha)\rceil$, which has probability
$k/(n+1)\ge1 - \alpha$.
\end{proof}

Exchangeability is exactly what a regime change breaks (\cref{fig:ml:uncert:coverage}). Around the boosted trees' mean
forecast, split conformal on the calibration year covers 90.7\% of the test days before the break and 74.0\% after it,
with a fixed width. Dividing the score by the ensemble's predicted standard deviation lets the width follow volatility:
90.3\% before, 86.9\% after. Adaptive conformal on the same normalised scores, updated day by day through the test
period with $\gamma = 0.005$, covers 89.9\% before and 90.1\% after (\cref{fig:ml:uncert:band}): it gives up the
finite-sample guarantee for a long-run one that survives drift.

\begin{omfigure}
\begin{tikzpicture}
\begin{axis}[width=12cm,height=5.4cm,ybar,bar width=7pt,ylabel={coverage of the 90\% interval (\%)},
  symbolic x coords={boosted quantiles,ensemble of 5,mixture density,conformal raw,conformal normalised,conformal adaptive},
  xtick=data,xticklabels={{boosted\\quantiles},{ensemble\\of five},{mixture\\density},{split\\conformal},{normalised\\conformal},{adaptive\\conformal}},
  xticklabel style={align=center,font=\footnotesize},tick label style={font=\small},label style={font=\small},
  ymin=70,ymax=95,grid=major,grid style={gray!25},enlarge x limits=0.1,legend columns=2,
  legend style={font=\footnotesize,at={(0.5,-0.3)},anchor=north,draw=none,fill=none,
  /tikz/every even column/.append style={column sep=8pt}},table/col sep=comma]
\addplot[fill=omDef!55,draw=omDef] table[x=name,y=before]{figdata/ml/11-probabilistic-models-and-uncertainty/coverage.csv};
\addplot[fill=omThm!55,draw=omThm] table[x=name,y=after]{figdata/ml/11-probabilistic-models-and-uncertainty/coverage.csv};
\draw[black!70,thick,dashed] ([xshift=-20pt]axis cs:boosted quantiles,90) -- ([xshift=20pt]axis cs:conformal adaptive,90);
\legend{before the regime change,after}
\end{axis}
\end{tikzpicture}

\omcaption{Coverage of 90\% intervals on the test days, before and after the volatility regime (dashed: the nominal 90\%).
The three models' own intervals, and conformal intervals around the boosted mean forecast: split conformal with raw
errors, with errors normalised by the ensemble's standard deviation, and adaptive. Data: \texttt{ml\_uncert.scores} and
\texttt{ml\_uncert.conformal}.}\label{fig:ml:uncert:coverage}
\end{omfigure}

\begin{omfigure}
\begin{tikzpicture}
\begin{axis}[width=12cm,height=5.4cm,xlabel={day},ylabel={next-day return (\%)},tick label style={font=\small},
  label style={font=\small},xmin=2000,xmax=2250,ymin=-7,ymax=7,xtick={2000,2050,2100,2150,2200,2250},grid=major,grid style={gray!25},
  xticklabel style={/pgf/number format/1000 sep={}},legend columns=3,legend style={font=\footnotesize,at={(0.5,-0.3)},
  anchor=north,draw=none,fill=none,/tikz/every even column/.append style={column sep=8pt}},table/col sep=comma]
\addplot[only marks,mark=*,mark size=0.9pt,black!70] table[x=day,y=y]{figdata/ml/11-probabilistic-models-and-uncertainty/band.csv};
\addplot[omDef,thick] table[x=day,y=hi]{figdata/ml/11-probabilistic-models-and-uncertainty/band.csv};
\addplot[omThm,thick,dashed] table[x=day,y=ahi]{figdata/ml/11-probabilistic-models-and-uncertainty/band.csv};
\addplot[omDef,thick,forget plot] table[x=day,y=lo]{figdata/ml/11-probabilistic-models-and-uncertainty/band.csv};
\addplot[omThm,thick,dashed,forget plot] table[x=day,y=alo]{figdata/ml/11-probabilistic-models-and-uncertainty/band.csv};
\draw[gray,thick] (axis cs:2100,-7) -- (axis cs:2100,7);
\legend{realised return,ensemble 90\% interval,adaptive conformal}
\end{axis}
\end{tikzpicture}

\omcaption{One asset through the regime change at day 2\,100 (vertical line): realised next-day returns, the ensemble's 90\%
interval, and the adaptive conformal interval around the boosted mean. Both widen after the break; the conformal band
widens further as its errors accumulate. Data: \texttt{ml\_uncert.conformal}.}\label{fig:ml:uncert:band}
\end{omfigure}

\section{Using uncertainty in sizing}

For independent bets, the weight that maximises the Sharpe ratio of their sum is proportional to the expected return over
the variance; Kelly sizing (Book~2, chapter~29) has the same form. The variance is where the model's distribution earns
its keep. The chapter's sizing test must be longer than its three test years, because Sharpe ratios estimated on three
years differ by more than the effect: on twenty-year series of twenty assets (four seeds, scored after the first seven
years), a daily portfolio weighted by the drift reading alone earns an annualised Sharpe ratio of 1.72 on average;
dividing the reading by a variance forecast raises it to 1.90 with the EWMA variance, 1.91 with the Gaussian network's
predicted variance and 1.93 with the true variance. With the true mean the gain is larger (2.94 to 3.31): the better the
mean forecast, the more the variance matters.

\begin{method}[Forecasting and using a distribution]\label{met:ml:uncert:method}
\begin{enumerate}
\item Train with a proper score (negative log-likelihood, pinball losses) and evaluate with another (CRPS) and with
coverage, on held-out periods including a stressed one.
\item Use an ensemble to separate epistemic from aleatoric variance; report both.
\item Wrap the forecast in conformal intervals, normalised by the predicted scale, and update them online.
\item Size by the mean over the variance, and test the rule on samples long enough for Sharpe ratios to differ.
\end{enumerate}
\end{method}

\section{Tutorial: sizing by doubt}\label{tut:ml:probabilistic-models-and-uncertainty:tut}

\begin{tutorial}
\textbf{Goal.} Fit quantile, Gaussian and mixture models of the next day's return, score them before and after a
volatility regime, build split and adaptive conformal intervals, and size positions by the predicted variance.
\textbf{End state:} \cref{tab:ml:uncert:scores}, \cref{fig:ml:uncert:coverage,fig:ml:uncert:band}, the sizing numbers.
\begin{tutsteps}
\item \textbf{Scores}: pinball, CRPS on a quantile grid, coverage.
\omcode{code/firm/uncert/firm_uncert.py}{32}{50}{The pinball loss, the CRPS and coverage.}\label{lst:ml:uncert:scores}
\item \textbf{Conformal}: the split threshold and the adaptive online level.
\omcode{code/firm/uncert/firm_uncert.py}{164}{185}{Split and adaptive conformal prediction.}\label{lst:ml:uncert:conformal}
\item \textbf{Run} \texttt{ml\_uncert.scores()}, \texttt{conformal()}, \texttt{sizing()} and \texttt{fig\_uncert.py}.
\end{tutsteps}
\textbf{What to change next.} Raise $\gamma$ to 0.02 and watch adaptive coverage recover faster and wobble more; size by
the mean over the standard deviation instead of the variance.
\end{tutorial}

\section{Build: predictive distributions and intervals}\label{bld:ml:probabilistic-models-and-uncertainty:uncert}

\begin{build}
\bfield{Purpose} Every forecast the firm uses comes with a distribution, a score for it, and an interval that holds up.
\bfield{Interface} \texttt{pinball}, \texttt{crps\_gaussian}, \texttt{crps\_quantiles}, \texttt{coverage};
\texttt{GaussNet}, \texttt{MDN}, \texttt{fit\_nll}, \texttt{predict\_gauss}, \texttt{predict\_mdn}, \texttt{mdn\_quantile};
\texttt{ensemble\_split}; \texttt{split\_conformal(scores, alpha)}, \texttt{adaptive\_conformal(scores, alpha, gamma,
window, start)}.
\bfield{Rules} Calibration data disjoint from training data and before the test; adaptive thresholds use only past
scores; every network deterministic on one thread.
\bfield{Acceptance tests} \texttt{code/firm/uncert/tests/}: the \omterm{def:ml:probabilistic-models-and-uncertainty:crps}{pinball loss} is minimised at the empirical quantile; the
Gaussian CRPS matches the quantile-grid approximation; split conformal covers at least $1 - \alpha$ on exchangeable data
across many draws; adaptive conformal restores coverage after a planted scale change; the Gaussian network and the mixture
recover a planted heteroskedastic and a bimodal distribution.
\bfield{Stretch} Conformalised \omterm{def:ml:probabilistic-models-and-uncertainty:quantile}{quantile regression}; conformal sets for classifiers (chapter~2's labels); a sizing rule
that shrinks the variance forecast.
\end{build}

\begin{omsources}
\item T.~Gneiting and A.~E.~Raftery, ``Strictly proper scoring rules, prediction, and estimation'', \emph{Journal of the
American Statistical Association} 102(477), 2007.
\item R.~Koenker and G.~Bassett, ``Regression quantiles'', \emph{Econometrica} 46(1), 1978.
\item C.~M.~Bishop, ``Mixture density networks'', Aston University technical report NCRG/94/004, 1994.
\item V.~Vovk, A.~Gammerman and G.~Shafer, \emph{Algorithmic Learning in a Random World}, Springer, 2005; 2nd ed., 2022.
\item I.~Gibbs and E.~Candès, ``Adaptive conformal inference under distribution shift'', \emph{NeurIPS}, 2021.
\item A.~Kendall and Y.~Gal, ``What uncertainties do we need in Bayesian deep learning for computer vision?'',
\emph{NeurIPS}, 2017.
\end{omsources}

\section{Exercises}

\begin{exercise}[$\star$]\label{exo:ml:probabilistic-models-and-uncertainty:1}
A 95\% quantile forecast is 2.0\%; the return is 3.0\%. What is the \omterm{def:ml:probabilistic-models-and-uncertainty:crps}{pinball loss}? And if the return is 1.0\%?
\end{exercise}

\begin{exercise}[$\star$]\label{exo:ml:probabilistic-models-and-uncertainty:2}
With 250 calibration scores and $\alpha = 0.1$, which order statistic is the split conformal threshold? With 5\,040?
\end{exercise}

\begin{exercise}[$\star$]\label{exo:ml:probabilistic-models-and-uncertainty:3}
Two bets have expected returns of 5 basis points and standard deviations of 80 and 250 basis points. In what proportion
should mean-variance sizing hold them?
\end{exercise}

\begin{exercise}[$\star\star$]\label{exo:ml:probabilistic-models-and-uncertainty:4}
Why does the raw split conformal interval lose so much coverage after the break, and the normalised one less?
\end{exercise}

\begin{exercise}[$\star\star$]\label{exo:ml:probabilistic-models-and-uncertainty:5}
The epistemic part of the predictive variance is 0.09\%. What does that say about the value of more data for this
forecast, and where would you expect it to be larger?
\end{exercise}

\begin{exercise}[$\star\star$]\label{exo:ml:probabilistic-models-and-uncertainty:6}
\emph{Find the flaw.} ``Sizing by the predicted variance lowered our Sharpe ratio from 2.1 to 1.8 over our three-year
test, so variance scaling does not work for this signal.''
\end{exercise}

\begin{exercise}[$\star\star\star$]\label{exo:ml:probabilistic-models-and-uncertainty:7}
\emph{Coding.} Run the adaptive conformal intervals of \texttt{ml\_uncert.conformal} with $\gamma = 0.02$. Report the
coverage after the break and explain the trade-off.
\end{exercise}

\begin{exercise}[$\star\star\star$]\label{exo:ml:probabilistic-models-and-uncertainty:8}
Show that the CRPS equals twice the integral over $\alpha$ of the \omterm{def:ml:probabilistic-models-and-uncertainty:crps}{pinball loss} of the forecast's $\alpha$-quantile.
\end{exercise}

\section{Problem: Sizing by Doubt}

\begin{problem}[{Weekend problem --- distributions through a regime change}]\label{pb:ml:probabilistic-models-and-uncertainty:1}
The chapter's twenty assets, the volatility regime from day 2\,100, and the four ways of forecasting the next day.

\medskip\noindent\textbf{Part I --- Scores.}
\begin{enumerate}
\item Why must a score be proper, and which of the chapter's scores are?
\item What do the models score (CRPS and coverage) before the break?
\item And after it?
\item Why are the CRPS so close across models and the truth?
\end{enumerate}

\medskip\noindent\textbf{Part II --- Uncertainty.}
\begin{enumerate}[resume]
\item What share of the predictive variance is epistemic?
\item What does a \omterm{def:ml:neural-networks-for-noisy-tabular-data:ensemble}{deep ensemble} add to a single Gaussian network here?
\item What does the \omterm{def:ml:probabilistic-models-and-uncertainty:quantile}{mixture density network} model that the Gaussian network does not?
\item Where would you expect the epistemic share to matter?
\end{enumerate}

\medskip\noindent\textbf{Part III --- Conformal.}
\begin{enumerate}[resume]
\item What do split conformal intervals cover before and after the break, raw and normalised?
\item What does adaptive conformal cover, and what does it give up?
\item Which assumption does the break violate?
\item How would you choose $\gamma$?
\end{enumerate}

\medskip\noindent\textbf{Part IV --- Sizing and the verdict.}
\begin{enumerate}[resume]
\item Why is the sizing test run on twenty-year series?
\item What Sharpe ratios do the four rules earn with the drift reading, and with the true mean?
\item Why does the variance matter more with a better mean forecast?
\item State the \emph{named result}: the Sharpe gain of variance-scaled sizing over sizing by the mean alone, and conformal
coverage against its nominal 90\% after the regime change.
\item What should a risk report show beside a point forecast?
\item How would you monitor a model's intervals in production (chapter~27)?
\item What would change with a fat-tailed target and a Gaussian model?
\item In one sentence: what is a forecast of a return?
\end{enumerate}
\end{problem}

\section{Interview questions}

\begin{interviewq}[$\star$ \iqroles{researcher, mle}]\label{iq:ml:probabilistic-models-and-uncertainty:1}
What is a \omterm{def:ml:probabilistic-models-and-uncertainty:scoring}{proper scoring rule}, and why should you train and evaluate forecasts with one?
\end{interviewq}

\begin{interviewq}[$\star\star$ \iqroles{mle}]\label{iq:ml:probabilistic-models-and-uncertainty:2}
How would you get prediction intervals out of a gradient-boosted model?
\end{interviewq}

\begin{interviewq}[$\star\star$ \iqroles{researcher}]\label{iq:ml:probabilistic-models-and-uncertainty:3}
Explain split \omterm{def:ml:probabilistic-models-and-uncertainty:conformal}{conformal prediction} and its guarantee. When does the guarantee fail in trading?
\end{interviewq}

\begin{interviewq}[$\star\star$ \iqroles{researcher, trader}]\label{iq:ml:probabilistic-models-and-uncertainty:4}
Your model forecasts both the mean and the variance of tomorrow's return. How do you size positions, and how would you
test the rule?
\end{interviewq}

\begin{interviewq}[$\star\star$ \iqroles{mle, researcher}]\label{iq:ml:probabilistic-models-and-uncertainty:5}
What is the difference between aleatoric and \omterm{def:ml:probabilistic-models-and-uncertainty:aleatoric}{epistemic uncertainty}, and how do you estimate each?
\end{interviewq}

\begin{interviewq}[$\star\star\star$ \iqroles{researcher}]\label{iq:ml:probabilistic-models-and-uncertainty:6}
Prove that the expected \omterm{def:ml:probabilistic-models-and-uncertainty:crps}{pinball loss} is minimised by the true quantile.
\end{interviewq}
